(1) 设管轴 \(x=0\) 为封闭端、\(x=L\) 为开口端。封闭端空气速度须为零，对线性声波意味着压力梯度 \(\partial_xp'(0,t)=0\)；理想开口相对外界压力扰动为零，即 \(p'(L,t)=0\)。因而压力空间振型为 \(P(x)\propto\cos(kx)\)，并要求 \(\cos(kL)=0\)。这里压力结点和速度结点不可混淆：它们在两端角色正好互换。
(2) 由 \(kL=(n-\tfrac12)\pi\)，\(n=1,2,\ldots\)，得 \(f_n=(2n-1)c/(4L)\)，其中基频 \(f_1=c/(4L)\)。频率相对于基频依次是 \(1,3,5,\ldots\)，理想开闭管没有偶数倍的纵向共振；这一谱结构来自两端不同的边界条件，不是介质本身产生非线性。
(3) 实际开口的空气运动略超出几何管口，需用 \(L_{\rm eff}=L+\delta\) 表示惯性延长。于是 \(f_n\simeq(2n-1)c/[4(L+0.61a)]\)，在 \(a/L\ll1\) 时为 \(f_n^{(0)}[1-0.61a/L+O((a/L)^2)]\)。经验系数约 \(0.61\) 仅用于无外法兰、细圆管且低频条件；不同口形或大 \(ka\) 时不能照搬常数。
(4) 定义 \(Q=\omega_n E/P_{\rm loss}\)（弱阻尼、周期平均意义），所以能量按 \(E(t)=E(0)e^{-t/\tau_E}\) 衰减，\(\tau_E=Q/\omega_n\)；由于能量与振幅平方成正比，振幅时间常数 \(\tau_A=2Q/\omega_n\)。在 \(Q\gg1\) 的洛伦兹响应中半功率频率全宽 \(\Delta\omega=\omega_n/Q\)，即 \(\Delta f=f_n/Q\)；三者用的是不同的衰减量，不能混为一谈。
